\section{Related Work}

Real-time fault diagnosis for cloud microservice systems has attracted increasing attention from both academia and industry. Existing studies can be broadly grouped into three categories, including rule-based and statistical methods, supervised deep learning methods, and recent large model-based diagnosis method. We review these lines of research and discuss their limitations in latency-sensitive cloud environments.

\subsection{Rule-Based and Statistical Methods}

Early studies mainly relied on handcrafted rules, statistical patterns, or explicit dependency structures for anomaly detection and troubleshooting. Xu et al. proposed a pioneering framework that mines console logs for detecting runtime system problems through feature extraction and machine learning analysis \cite{xu2009detecting}. Lou et al. further introduced invariant mining from console logs to identify abnormal executions in distributed systems \cite{lou2010mining}. Lin et al. developed LogCluster, which clusters representative log sequences and leverages historical knowledge bases to support problem identification in online services \cite{lin2016log}. In the context of microservice systems, Brandón et al. proposed a graph-based root cause analysis framework that matches anomalous dependency graphs against a predefined anomaly library \cite{brandon2020graph}. More recently, BARO integrated multivariate Bayesian online change-point detection with statistical hypothesis testing for robust troubleshooting in microservices \cite{pham2024baro}. Statistical prediction techniques have also been explored in datacenter infrastructure management, such as PreFix for switch failure prediction using syslog patterns and temporal statistics \cite{zhang2018prefix}. Comprehensive empirical evaluations of traditional log anomaly detection methods were reported in \cite{he2016experience}.

These methods are generally lightweight, interpretable, and easy to deploy. However, they often depend on handcrafted thresholds, manually curated rules, stationary assumptions, or predefined causal templates. Such dependencies limit their scalability and maintainability in large, dynamic, and continuously evolving cloud microservice systems.

\subsection{Supervised Deep Learning Methods}

To reduce manual feature engineering and improve diagnostic accuracy, recent works increasingly employ deep learning for anomaly detection and root cause localization. DeepLog introduced an LSTM-based framework that models log sequences as language-like data and performs online anomaly detection and workflow diagnosis \cite{du2017deeplog}. LogAnomaly further incorporated semantic embeddings and quantitative patterns to jointly detect sequential and numerical anomalies in logs \cite{meng2019loganomaly}. For microservice troubleshooting, Eadro proposed the first end-to-end framework that jointly performs anomaly detection and root-cause localization using traces, logs, and KPIs via multi-task learning \cite{10172617}. Beyond cloud systems, deep models have also been widely studied in industrial fault diagnosis and intelligent monitoring. Representative examples include real-time marine machinery diagnosis using variational autoencoders \cite{VELASCOGALLEGO2022117634}, anomaly-aware autonomous vehicle diagnosis with denoising autoencoders \cite{MIN2023120002}, diffusion-based wind turbine fault identification with SHAP explanations \cite{YAO2025127925}, lightweight hardware-aware bearing anomaly diagnosis \cite{11040080}, hydro-turbine acoustic anomaly detection with fusion networks \cite{10858663}, semi-supervised anomaly detection via restricted distribution transformation \cite{11071974}, spatiotemporal anomaly detection and localization for distributed parameter systems \cite{10829379}, and large-scale offshore wind turbine diagnosis using Informer and XGBoost \cite{10978904}.

Although deep learning methods significantly improve diagnosis effectiveness through automatic representation learning, many approaches rely on heavy end-to-end architectures whose inference overhead grows rapidly with data scale, system size, or modality complexity. This limits their practicality for latency-sensitive online diagnosis in production cloud environments.

\subsection{Large Model-Based Diagnosis Method}

Recently, foundation models and large language models have opened new opportunities for intelligent system operations due to their strong semantic understanding, transfer learning capability, and reasoning capacity. LogLLM employed pre-trained language models for log anomaly detection by aligning semantic embeddings with autoregressive sequence classification, demonstrating strong robustness on unstable logs \cite{guan2024logllm}. ICAD-LLM introduced an in-context anomaly detection paradigm that uses large language models to unify anomaly detection across heterogeneous modalities and unseen domains \cite{wu2026icad}. Beyond log intelligence, multimodal large models have also been explored for time-series anomaly detection. Recent studies investigated whether multimodal LLMs can detect point-wise, range-wise, and irregular anomalies through visual reasoning, and proposed agent-based frameworks that combine LLM reasoning with external tools for automatic diagnosis \cite{xu2026can}. In software reliability analysis, graph-enhanced semantic models such as Lograph have further demonstrated the benefit of combining structured system entities with semantic understanding for anomaly detection \cite{10839428}.

Despite their promising reasoning ability and cross-domain generalization, foundation-model-based methods often incur substantial computational cost, memory overhead, and deployment complexity. These issues make direct adoption challenging in latency-sensitive cloud environments where bounded latency, efficient resource usage, and runtime responsiveness are critical.

\subsection{Summary and Motivation}

In summary, rule-based and statistical methods are efficient but brittle under evolving environments. Deep learning methods improve accuracy but often suffer from high inference cost. Foundation-model-based approaches offer stronger intelligence but remain expensive for online deployment. Therefore, there is still a clear need for a practical diagnosis framework that jointly achieves multimodal effectiveness, low-latency inference, adaptive computation, and deployment efficiency for modern cloud microservice systems. Motivated by this gap, we propose an online multimodal fault diagnosis framework via adaptive sparse inference.
